\section{Proof of Theorem~\ref{thm: extmain}}\label{sec: main}
In this section, we prove Theorem~\ref{thm: extmain}.  We first formally define the function we use.

\subsection{Definition of the function}\label{sec: Fdefn}
If $\delta < 1/2$, then ignore the last $n - 2n^{2\delta}$ bits of the input, and define the following function on the first $2n^{\delta}$ bits of the input. The same argument as in Sections~\ref{sec: upper} and~\ref{sec: lower} give the required bounds for Theorem~\ref{thm: extmain} and Theorem~\ref{thm: wtdeg}.
Hence, we may assume without loss of generality that $\delta \geq 1/2$.

Define the partial function $f : \pmone^n \to \bra{-1, 1, \star}$ by $f = \PARITY_{n^{\delta}/2}^{\HADD_{n^{1 - \delta}}}$.
Define $F$ to be the completion of $f$ that evaluates to $-1$ on the non-promise domain of $f$ (see Figure~\ref{fig: cexample}).




\begin{figure}[htbp]
\begin{center}
\begin{tikzpicture}

\tikzstyle{gate}=[ellipse,draw=black]
\tikzstyle{input}=[]
\node(hdef) at (-6, -1){$F = $};
\node[gate] (root) at (0,0) {$\mathsf{PARITY}$};
\node[gate] (c1) at (-3,-2) {$\mathsf{HADD}_\ell$};
\node[gate] (cn) at (3,-2) {$\HADD_{\ell}$};
\node (dummy1) at (-1.8, -0.3) {};
\node (dummy2) at (1.9, -0.3) {$k/2$};
\draw[->, dashed] (dummy1) edge [bend right = 20] (dummy2);

\node (dummy3) at (-5, -2.7) {$2\ell$};
\node (dummy4) at (-1, -2.8) {};
\draw[<-, dashed] (dummy3) edge [bend right = 20] (dummy4);


\node[input] (x8) at (1,-2) {$\boldsymbol{\cdot}$};
\node[input] (x19) at (-1,-2) {$\boldsymbol{\cdot}$};
\node[input] (x419) at (0,-2) {$\boldsymbol{\cdot}$};

\node[input] (x11) at (-5.5,-4) {$x_{11}$};
\node[input] (x123) at (-4.5,-4) {$\cdots$};
\node[input] (x1f) at (-3.5,-4) {$x_{1\ell}$};
\node[input] (y11) at (-2.5,-4) {$y_{11}$};
\node[input] (y123) at (-1.5,-4) {$\cdots$};
\node[input] (y1f) at (-0.5,-4) {$y_{1\ell}$};


\node[input] (yn1) at (5.5,-4) {$y_{\frac{k}{2}\ell}$};
\node[input] (x13) at (4.5,-4) {$\cdots$};
\node[input] (ynf) at (3.5,-4) {$y_{\frac{k}{2}1}$};
\node[input] (xn1) at (2.5,-4) {$x_{\frac{k}{2}\ell}$};
\node[input] (y13) at (1.5,-4) {$\cdots$};
\node[input] (xnf) at (0.5,-4) {$x_{\frac{k}{2}1}$};

\draw[->] (c1) -- (root);
\draw[->] (cn) -- (root);


\draw[<-] (c1) -- (y11);
\draw[<-] (c1) -- (y1f);
\draw[<-] (c1) -- (x11);
\draw[<-] (c1) -- (x1f);

\draw[<-] (cn) -- (yn1);
\draw[<-] (cn) -- (ynf);
\draw[<-] (cn) -- (xn1);
\draw[<-] (cn) -- (xnf);

\draw[->] (root) -- ++(0,1);

\draw [
    thick,
    decoration={
        brace,
        mirror,
        raise=0.5cm
    },
    decorate
] (-5.8,-4) -- (-3.2, -4)
node [pos=0.5,anchor=north,yshift=-0.55cm] {Address bits}; 
\draw [
    thick,
    decoration={
        brace,
        mirror,
        raise=0.5cm
    },
    decorate
] (-2.8,-4) -- (-0.3, -4)
node [pos=0.5,anchor=north,yshift=-0.55cm] {Target bits}; 
\end{tikzpicture}
\end{center}
\caption{$k = n^{\delta}, \ell = n^{1 - \delta}$. If the address bits of an input to an $\HADD_\ell$ is the $j$'th Hadamard codeword, then $y_j$ is selected.  If on an input, there exists at least one $\HADD_\ell$ for which the address bits do not correspond to a Hadamard codeword, $F$ outputs $-1$.  Else it outputs the Parity of the $k/2$ selected target bits.}
\label{fig: cexample}
\end{figure}



\subsection{Upper bound}\label{sec: upper}
In this section, we prove the following.
\begin{claim}\label{claim: ndelta}
For $F: \pmone^{n} \to \pmone$ defined as in Section~\ref{sec: Fdefn}, we have
\[
Q(F) = \Theta(n^{\delta}).
\]
\end{claim}
The upper bound follows along the lines of a proof in~\cite{AdW14}, and the lower bound just uses the fact that $F$ is at least as hard as $\PARITY_{n^{\delta}/2}$.
\begin{proof}
For convenience, set $\ell = n^{1-\delta}$ and $k = n^{\delta}$.
Recall that the input to $F$ is viewed as $(x_{11},\ldots,x_{1\ell},y_{11},\ldots,y_{1\ell},\ldots,x_{\frac{k}{2}1},\ldots,x_{\frac{k}{2}\ell},y_{\frac{k}{2}1},\ldots,y_{\frac{k}{2}\ell})$. The following is an $O(n^\delta)$-query quantum algorithm computing $F$. Note that since $\delta \geq 1/2$, we have $k = \Omega(\ell)$.
\begin{enumerate}
    \item Run $k/2$ instances of Bernstein-Vazirani algorithm on inputs $(x_{11}, \dots, x_{1\ell}), \dots, (x_{\frac{k}{2}1}, \dots, x_{\frac{k}{2}\ell})$ to obtain $k/2$ strings $z_1, \dots, z_{k/2}$.
    \item Run Grover's search~\cite{Grover96, BHMT02} to check equality of the two strings: $h(z_1), \dots, h(z_\ell)$ and  $x_{11}, \dots, x_{1\ell}, \dots, x_{\frac{k}{2}1}, \dots, x_{\frac{k}{2}\ell}$, i.e.~to check whether the addressing bits of the input are indeed all Hadamard codewords which are output by the first step.
    \item If the step above outputs that the strings are equal, then query the $k/2$ selected variables and output their parity.  Else, output $-1$.
    \begin{itemize}
        \item If the input was indeed of the form as claimed in the first step, then Bernstein-Vazirani outputs the correct $z_1, \dots, z_\ell$ with probability 1, and Grover's search verifies that the strings are equal with probability 1.  Hence the algorithm is correct with probability 1 in this case.
        \item If the input was not of the claimed form, then the two strings for which equality is to be checked in the second step are not equal.  Grover's search catches a discrepancy with probability at least $2/3$.  Hence, the algorithm is correct with probability at least $2/3$ in this case.
        \end{itemize}
\end{enumerate}
The correctness of the algorithm is argued above, and the cost is $k/2$ queries for the first step, $O(\sqrt{k\ell})$ queries for the second step, and at most $k/2$ for the third step.
Thus, we have
\[
Q(F) = O(k + \sqrt{k\ell}) = O(k),
\]
since $k = \Omega(\ell)$.  The upper bound in the lemma follows.

For the lower bound, we argue that $F$ is at least as hard as Parity on $k/2$ inputs. To see this formally, set all the address variables such that the selected target variables are the first target variable in each block.  Under this restriction, $F$ equals $\mathsf{PARITY}(y_{11}, \dots, y_{\frac{k}{2}1})$.  Thus any quantum query algorithm computing $F$ must be able to compute $\mathsf{PARITY}_{k/2}$, and thus $Q(F)=\Omega(k)$.
\end{proof}

\begin{remark}\label{rmk: outer}
The same argument as above works when the function $f$ is defined to be $g^{\HADD_\ell}$ for any $g : \pmone^{n^{\delta}} \to \pmone$ satisfying $\adeg(g) = \Omega(n^{\delta})$, and $F$ is the completion of $f$ that evaluates to $-1$ on all non-promise inputs.  The same proof of Theorem~\ref{thm: extmain} also goes through, but we fix $g = \PARITY_{n^{\delta}/2}$ for convenience.
\end{remark}

\subsection{Lower bound}\label{sec: lower}

In this section, we first prove Lemma~\ref{lem: lifting}.
We require the following observation.
\begin{observation}\label{obs: obs}
For any monomial $\chi_S,~S\subseteq [n]$ and any $j\in S$, we have 
\[
\Exp_{x_j \sim \pmone}[\chi_S(x)] = 0,
\]
where $x_j$ is distributed uniformly over $\pmone$.
\end{observation}

\begin{proof}[Proof of Lemma~\ref{lem: lifting}]

Let $F = f^{\ADD_{m, t}}$. Recall that our goal is to show that $\log \asnorm{F}{1/3} = \Omega(\adeg(f) \log t)$. We may assume $\adeg(f) \geq 1$, because the lemma is trivially true otherwise.  

Towards a contradiction, suppose there exists a polynomial $P$ of spectral norm strictly less than $2^{\frac{1}{10}\adeg_{0.99}(f)\log t}$ uniformly approximating $F$ to error $1/3$ on the promise inputs (recall that from Lemma~\ref{lem: asnormequiv}, we have $\adeg(f) = \Theta(\adeg_{0.99}(f))$).

Let $\nu$ be a distribution on the address bits of $\ADD_{m, t}$ such that $\nu$ is supported only on assignments to the address variables that do not select $\star$, and is the uniform distribution over these assignments.
Let $\mu = \nu^n$ be the product distribution over the address bits of the addressing functions in $F$.

\begin{itemize}
    \item For any assignment $z$ of the address variables from the support of $\mu$, define a relevant (target) variable to be one that is selected by $z$. Analogously, define a target variable to be irrelevant if it is not selected by $z$. Define a monomial to be relevant if it does not contain irrelevant variables, and irrelevant otherwise.
    \item Note that for any target variable, the probability with which it is selected is exactly $1/t$.
    \item We are interested in the set of target variables present in a monomial $\chi_S$. We denote this set by $S_{\mathrm{target}}$. Define the target-degree of $S$ to be $|S_{\mathrm{target}}|$, the degree when restricted to the target variables.
    \item Thus under any assignment $z$ drawn from $\mu$, for any monomial of $P$ of target-degree $t \geq \adeg_{0.99}(f)$, the probability that it is relevant is at most $1/t^{\adeg_{0.99}(f)}$.
    \item Hence, 
    \begin{align*}
    & \Exp\limits_{z\sim\mu}[\begin{gathered}\text{$\ell_1$-norm of relevant monomials in }P\\\text{of target-degree }\geq\adeg_{0.99}(f)\end{gathered}]\\
    &= \sum\limits_{|S_{\mathrm{target}}| \geq \adeg_{0.99}(f)} |\widehat P(S)| \Pr\limits_{z\sim\mu}[\chi_S \text{ is relevant w.r.t. }z] \\
    &\leq \max_{|S_{\mathrm{target}}|\geq\adeg_{0.99}(f)}\{\Pr\limits_{z\sim\mu}[\chi_S \text{ is relevant w.r.t. }z]\} \cdot \snorm{P} \\
    &< \frac{1}{t^{\adeg_{0.99}(f)}} \cdot 2^{\frac{1}{10}\adeg_{0.99}(f) \log t} \\
    & = 2^{(-\frac{9}{10})\adeg_{0.99}(f)\log t}< \frac{3}{5},
    \end{align*}
    where the second last inequality holds from Claim~\ref{claim: adegnaivelb}, and the last inequality holds because $t\geq 2$ and $\adeg_{0.99}(f) \geq 1$.
    \item Fix an assignment to the address variables from the support of $\mu$ such that under this assignment, the $\ell_1$-norm of the relevant monomials in $P$ of degree $\geq \adeg_{0.99}(f)$ is less than $3/5$.
    \item Note that under this assignment (in fact under any assignment in the support of $\mu$), the restricted $F$ is just the function $f$ on the $n$ variables selected by the addressing functions. Denote by $P_1$ the polynomial on the target variables obtained from $P$ by fixing address variables as per this assignment.
    \item Drop the relevant monomials of degree $\geq \adeg_{0.99}(f)$ from $P_1$ to get a polynomial $P_2$, which uniformly approximates the restricted $F$ (which is $f$ on $n$ variables) to error $1/3 + 3/5 < 0.99$.
    \item Take expectation over irrelevant variables (from the distribution where each irrelevant variable independently takes values uniformly from $\pmone$).  Under this expectation, the value of $F$ does not change (since irrelevant variables do not affect $F$'s output by definition), and all irrelevant monomials of $P_2$ become 0 (using Observation~\ref{obs: obs} and linearity of expectation).  Hence, under this expectation we have $\Exp[P_2] = P_3$, where $P_3$ is a polynomial of degree strictly less than $\adeg_{0.99}(f)$.  Furthermore, $P_3$ uniformly approximates $f$ to error less than $0.99$ which is a contradiction. 
\end{itemize}
\end{proof}


As a corollary of Lemma~\ref{lem: lifting}, we obtain a lower bound on the approximate spectral norm of $F$, where $F$ is defined as in Section~\ref{sec: Fdefn}.  This yields a proof of Theorem~\ref{thm: wtdeg}.

\begin{proof}[Proof of Theorem~\ref{thm: wtdeg}]
Given $0 < \delta < 1$, construct $F$ as in Section~\ref{sec: Fdefn}.  Claim~\ref{claim: ndelta} implies 
\[
Q(F) = \Theta(n^{\delta}).
\]
Let $f = \PARITY_{n^{\delta}/2}^{\HADD_{n^{1 - \delta}}}$.  Lemma~\ref{lem: lifting} implies that 
\[
\asnorm{f}{1/3} = \Omega(n^{\delta} \log n).
\]
Since $F$ is a completion of $f$, we have 
\[
\asnorm{F}{1/3} = \Omega(n^{\delta} \log n),
\]
which proves the lower bound in Theorem~\ref{thm: wtdeg}.  The upper bound follows from Theorem~\ref{thm: qtm_sim_log_loss}.
\end{proof}

We are now ready to prove our main theorem.
\begin{proof}[Proof of Theorem~\ref{thm: extmain}]
It immediately follows from Theorem~\ref{thm: wtdeg} and Theorem~\ref{thm: LS}.
\end{proof}

